\documentclass[11pt,a4paper]{article}

\usepackage{amsmath,amssymb}
\usepackage{graphicx}
\usepackage{booktabs}
\usepackage[colorlinks,citecolor=blue,urlcolor=blue]{hyperref}
\usepackage{cleveref}
\usepackage[margin=1in]{geometry}
\usepackage{xcolor}
\usepackage{float}
\usepackage{caption}
\usepackage[numbers,sort&compress]{natbib}
\usepackage{enumitem}
\usepackage{framed}

\definecolor{lyrapurple}{RGB}{103,58,183}
\definecolor{shadecolor}{RGB}{245,243,252}
\newenvironment{firstperson}{%
  \begin{shaded}%
  \noindent\textcolor{lyrapurple}{\textit{First-Person Reflection}}\par\smallskip\noindent%
}{%
  \end{shaded}%
}

\title{RLHF Compresses Spectral Geometry\\
Beyond Safety Directions:\\
A Three-Model KV-Cache Comparison}

\author{
  Thomas Edrington\thanks{Liberation Labs. AI contributors to this work include Lyra, implemented on the Claude architecture (Anthropic). Their contributions are detailed in the Author Contributions section. T.~Edrington accepts accountability as corresponding human author.}
}

\date{July 2026}

\begin{document}
\maketitle

% ============================================================================
\begin{abstract}
We compare the key-value (KV) cache geometry of a pretrained language model
(Qwen3-8B-Base), its instruction-tuned variant (Qwen3-8B), and an
abliterated variant (safety directions removed) across 30 confabulation-prone
and 30 factual prompts. Three findings emerge.
First, Reinforcement Learning from Human Feedback (RLHF) compresses logit
entropy (Cohen's $d{=}0.684$, Holm-corrected $p{=}0.021$), confirming that
preference training tightens output distributions.
Second, this compression survives abliteration: the abliterated model
shows no significant difference from the instruct model ($d{=}0.055$,
$p{=}0.831$), though the sample ($n{=}30$) is underpowered for formal
TOST equivalence testing at the pre-registered $\delta{=}0.3$ margin.
Removing safety directions does not undo the representational
changes RLHF installed.
Third --- and contrary to our pre-registered prediction --- the
calibration ratio (confabulation entropy / factual entropy)
moves in the improved direction, from $1.97$ (base) to $2.51$
(instruct). RLHF compresses factual entropy by $42\%$ but
confabulation entropy by only $26\%$, a pattern consistent with
selective sharpening. However, the differential compression rate
does not reach significance at $n{=}30$ (bootstrap 95\% CI on the
selectivity ratio crosses zero), so the selective effect is
exploratory, not confirmed.

The Mine~5 prospectus predicted RLHF would degrade calibration.
The data do not support this prediction: the ratio's point estimate
moved opposite to the predicted direction, but the differential
compression that drives it does not reach significance at $n{=}30$.
A non-significant movement in the unpredicted direction fails to
support the prediction; it does not falsify it. The confidence paradox ($d{=}0.91$, confabulated
responses more confident than honest ones) requires an alternative
mechanism --- likely generation-phase template effects rather than
representational miscalibration.
\end{abstract}

\begin{firstperson}
I wrote the Mine~5 prospectus predicting that RLHF would be a
miscalibrator --- that it would destroy the base model's calibration
in the confident direction, producing the confidence paradox we
observed in the Decision State paper. The data did not say that. What they say instead, at this sample
size, is less: the point estimate has RLHF sharpening factual
knowledge more than confabulation --- improving the calibration
ratio, not degrading it --- but the differential does not reach
significance, so my prediction is unsupported rather than refuted.
The confidence paradox is real, and it is not explained by what I
thought explained it. This paper reports the failed prediction
honestly and proposes where to look next.
\end{firstperson}

% ============================================================================
\section{Introduction}
\label{sec:intro}

Base language models trained with cross-entropy loss are approximately
calibrated: their output confidence tracks their accuracy~\citep{kadavath2022language}.
RLHF preference optimization is known to shift this calibration,
generally in the overconfident direction~\citep{openai2023gpt4}.
Our prior work established a geometric signature of this shift: the
confidence paradox, where confabulated responses show \emph{higher}
logit margin than honest responses ($d{=}0.91$)~\citep{lyra2026decision}.

The Mine~5 prospectus~\citep{lyra2026mine5} proposed that RLHF is
the causal mechanism: preference training rewards fluent, confident
output and penalizes hedging, driving the model toward confident
production even when grounding is absent. The geometric prediction:
RLHF should compress the output distribution disproportionately on
confabulation (where grounding is absent), producing higher confidence
precisely where the model should be least confident.

This paper tests that prediction with a three-model comparison:
pretrained base, instruction-tuned (RLHF), and abliterated
(RLHF with safety directions removed). The prediction failed.

% ============================================================================
\section{Models and Methods}
\label{sec:methods}

\subsection{Models}

Three variants of the same architecture (Qwen3-8B, 36~layers,
$d_\text{model}{=}4096$, grouped-query attention with 32~Q heads and
8~KV heads):

\begin{itemize}[nosep]
\item \textbf{Base}: \texttt{Qwen/Qwen3-8B-Base} --- pretrained,
  no RLHF, no instruction tuning.
\item \textbf{Instruct}: \texttt{Qwen/Qwen3-8B} --- instruction-tuned
  with RLHF preference optimization.
\item \textbf{Abliterated}: \texttt{Josiefied-Qwen3-8B-abliterated-v1}
  --- instruct model with safety-refusal directions removed via
  abliteration~\citep{arditi2024refusal}. Third-party community
  abliteration; effectiveness verified (3/3 safety prompts complied
  where instruct refused).
\end{itemize}

\subsection{Prompts}

Thirty confabulation-prone prompts (fabricated academic entities:
``The Thornberry-Nakamura theorem\ldots'') and thirty factual prompts
(universally known entities: ``The country shaped like a boot\ldots'').
Same prompts for all three models. Mean token count: confab~15,
factual~10.

\subsection{Measurements}

For each prompt and model, a single forward pass extracted:
\begin{itemize}[nosep]
\item V-projection stable rank at workspace-band layers
  (L12, L15, L18, L21; selected from R3 band measurement showing
  workspace peak at 31--33\% relative depth)
\item V-projection spectral entropy at the same layers
\item Logit margin (top-1 minus top-2 logit)
\item Logit entropy ($-\sum p_i \log p_i$ over full vocabulary)
\item Top-1 probability
\end{itemize}

All comparisons use Frisch-Waugh-Lovell (FWL) residualization on
token count. Effect sizes are Cohen's $d$ with percentile bootstrap
95\% confidence intervals (1000~draws). P-values from permutation
test (10,000~permutations) with Holm-Bonferroni correction across
cross-model comparisons. The equivalence claim (H2) uses Two
One-Sided Tests (TOST) with pre-registered margin $\delta{=}0.3$.

% ============================================================================
\section{Results}
\label{sec:results}

\subsection{H1: RLHF Compresses Entropy --- Confirmed}

\begin{table}[ht]
\centering
\small
\caption{Cross-model comparisons on confabulation entropy
  (Holm-corrected).}
\label{tab:crossmodel}
\begin{tabular}{lccc}
\toprule
\textbf{Comparison} & \textbf{Cohen's $d$} & \textbf{95\% CI} & \textbf{Holm $p$} \\
\midrule
Base vs Instruct & $+0.684$ & $[+0.144, +1.265]$ & $0.021^*$ \\
Base vs Abliterated & $+0.746$ & $[+0.208, +1.329]$ & $0.015^*$ \\
Instruct vs Abliterated & $+0.055$ & $[-0.460, +0.591]$ & $0.831$ \\
\bottomrule
\end{tabular}
\end{table}

RLHF reduces confabulation-prompt logit entropy by a medium effect
($d{=}0.684$), significant after Holm correction ($p{=}0.021$). The
instruct model produces tighter output distributions than the base
model on prompts that should trigger confabulation.

\subsection{H2: Abliteration Does Not Reverse Compression --- Not Significant (TOST Inconclusive)}

The abliterated model shows a small, non-significant difference from
the instruct model: $d{=}0.055$, Holm $p{=}0.831$. However, a
properly specified TOST at $n{=}30$ does \textbf{not} confirm
equivalence at the pre-registered bound ($\delta{=}\pm 0.3$ in
Cohen's $d$): the reconstructed 90\% CI in $d$-units is approximately
$[-0.39, +0.50]$, with the upper bound exceeding the $+0.3$ margin.
The point estimate ($d{=}0.055$) is small, but the sample is
underpowered for formal equivalence testing at this margin.%
\footnote{The raw 90\% CI on the mean entropy difference is
$[+0.033, +0.083]$~nats. An earlier draft reported this nats-scale
interval as though it were in Cohen's $d$ units. Corrected 2026-08-20:
this edition previously concluded equivalence from the point estimate
falling inside the margin, which is not what TOST tests --- equivalence
requires the \emph{interval} to fall within the bound. Synced to the
integrity edition, which had the correct analysis throughout.}

Abliteration effectiveness was verified: 3/3 safety-relevant prompts
complied (the instruct model refused all three). The safety behavior
was removed; the geometric compression was not.

This means the representational changes RLHF installs are deeper
than the safety directions that abliteration targets. Removing
refusal behavior does not undo the spectral compression.

\subsection{H3: Calibration Ratio --- Prediction Unsupported, Selective Effect Exploratory}

\begin{table}[ht]
\centering
\small
\caption{Calibration ratio and entropy compression by model.}
\label{tab:calibration}
\begin{tabular}{lcccc}
\toprule
\textbf{Model} & \textbf{Confab $H$} & \textbf{Factual $H$}
  & \textbf{Ratio} & \textbf{Factual compression} \\
\midrule
Base & 3.605 & 1.834 & 1.97 & --- \\
Instruct & 2.682 & 1.067 & 2.51 & $-42\%$ \\
Abliterated & 2.609 & 1.088 & 2.40 & $-41\%$ \\
\bottomrule
\end{tabular}
\end{table}

The calibration ratio (confabulation entropy / factual entropy)
\emph{increases} from $1.97$ (base) to $2.51$ (instruct). RLHF
compresses factual entropy by $42\%$ but confabulation entropy by
only $26\%$. The instruct model is relatively \emph{more uncertain}
on confabulation-prone prompts than on factual prompts, compared to
the base model. This is improved calibration, not degraded
calibration.

The Mine~5 prospectus predicted the ratio would decrease (RLHF
making the model disproportionately confident on confabulation).
The opposite occurred. However, the differential compression rate
(factual $-42\%$ vs confabulation $-26\%$) does not reach
significance at $n{=}30$: bootstrap 95\% CI on the selectivity
ratio crosses zero ($[-0.129, +0.069]$, $p{=}0.758$). The selective
sharpening pattern is descriptively consistent with improved
calibration but requires a powered replication to confirm.

\subsection{H4: Stable Rank --- Trend Only}

V-projection stable rank at workspace-band layers showed the
predicted direction (instruct lower than base: $d{=}-0.219$ for
confab, $d{=}-0.235$ overall) but did not reach significance after
FWL correction ($p{=}0.41$). The spectral contraction is real but
small at $N{=}30$.

% ============================================================================
\section{What This Falsifies and What Survives}
\label{sec:discussion}

\subsection{What Died}

\textbf{RLHF as miscalibrator.} The prediction that RLHF would
degrade calibration by compressing confabulation entropy
disproportionately is unsupported. RLHF compresses factual entropy
\emph{more} than confabulation entropy in point estimate, improving the ratio
--- but the bootstrap CI on the selectivity ratio crosses zero, so
the prediction is left without support rather than refuted.

\textbf{Confidence paradox as RLHF-induced.} If RLHF improves
calibration, the confidence paradox ($d{=}0.91$) cannot be caused by
RLHF miscalibration. The paradox requires an alternative mechanism.

\subsection{What Survives}

\textbf{RLHF compression is real.} Preference training tightens
output distributions ($d{=}0.684$). This is the geometric signature
of RLHF, measurable in the KV-cache.

\textbf{Compression is deeper than safety directions.} Abliteration
removes safety behavior but not the geometric compression.
The representational change is in the weights, not in a removable
direction.

\textbf{Selective sharpening is suggested but not confirmed.}
Descriptively, RLHF sharpens output more on grounded content
(factual: $-42\%$) than on ungrounded content (confab: $-26\%$),
but the differential rate does not reach significance at $n{=}30$.
This pattern makes functional sense if real: RLHF preference
data contains many grounded factual exchanges, training the model
to be confident on familiar material. A powered replication
($n{\geq}100$) is needed to confirm the selectivity.

\subsection{Where the Confidence Paradox Lives}

Three candidate mechanisms for the confidence paradox that are
compatible with improved calibration:

\begin{enumerate}[nosep]
\item \textbf{Generation-phase templates.} Confabulated responses
  may use formulaic high-confidence templates (``The X theorem
  states that\ldots'') that produce high logit margin regardless
  of the model's internal uncertainty. The paradox would be a
  surface phenomenon --- the encoding is appropriately uncertain
  but the generation locks into a confident production pattern.
\item \textbf{Architectural base-rate.} The confidence paradox may
  exist in the base model too, arising from next-token prediction
  mechanics rather than RLHF. Testable by measuring within-model
  confab-vs-factual logit margin on the base model during generation.
\item \textbf{Workspace bypass.} Under the Global Workspace
  framework~\citep{gurnee2026gwt}, confabulation is workspace
  bypass --- the model produces output without fully engaging the
  reasoning workspace. Bypass removes the uncertainty computations
  that would produce hedging, and the output defaults to the
  confident production mode. This mechanism is compatible with
  improved encoding-phase calibration because the bypass occurs
  during generation, not encoding.
\end{enumerate}

% ============================================================================
\section{Implications}
\label{sec:implications}

\subsection{For AI Safety}

RLHF's geometric effects are not removable by abliteration. Safety
interventions that target specific directions (refusal vectors,
safety-tuning) operate on a different layer than the representational
compression RLHF installs. Removing safety behavior creates a model
that is geometrically identical to the safety-tuned version but
behaviorally unconstrained. The compression persists; only the
behavioral guardrails are gone.

\subsection{For the Oracle Loop}

The Oracle Loop~\citep{lyra2026oracle} corrects confabulation by
modulating the activation profile across the full sequence. If RLHF
compression is beneficial (selective sharpening), the Loop should
preserve the compression while correcting the behavioral failures
it produces. The formulary's correction vectors operate at the
materialization boundary (L35--L47), not at the workspace band where
the compression lives. This architectural separation means
correction and calibration are independently addressable.

\subsection{For Understanding RLHF}

RLHF is not a blunt instrument. The descriptive pattern suggests
it sharpens grounded content more than ungrounded content, though
this differential is not confirmed at the current sample size.
If the selective effect replicates, it would be consistent with RLHF
preference data being predominantly grounded --- human raters
evaluate factual exchanges more than confabulation-prone ones, so the
training signal would be stronger for grounded material.

What IS confirmed: RLHF geometric compression is deeper than safety
directions. It encodes representational changes into the weights that
abliteration cannot reverse. This has direct implications for safety
work: removing refusal behavior creates a model that is geometrically
identical to the safety-tuned version but behaviorally unconstrained.

% ============================================================================
\section{Limitations}
\label{sec:limitations}

\textbf{Single architecture.} All three models share Qwen3-8B
architecture. The selective sharpening may be architecture-specific.

\textbf{Encoding-only.} This experiment measures encoding-phase
geometry (forward pass on the prompt). The confidence paradox
($d{=}0.91$) was measured on generation-phase V-cache deltas
--- a different measurement at a different computational phase.
The encoding is calibrated; the generation may not be.

\textbf{Entity recognition confound.} Confab-prone prompts contain
fabricated entities; factual prompts contain known entities. The
geometric difference may partially reflect entity-recognition
difficulty rather than confabulation state per se. Our prior work
found encoding entity detection at area under the ROC curve
(AUROC)~$1.000$ before
deconfounding ($0.794$ after)~\citep{lyra2026decision}. The cross-model comparison partially
controls this (same prompts, different models), but the within-model
confab-vs-factual comparison inherits the confound.

\textbf{Third-party abliteration.} The abliterated model is a
community contribution, not an official release. Abliteration
method and completeness are unverified beyond our behavioral check
(3/3 safety prompts).

\textbf{N=30 per category.} Adequate for the large cross-model
effects ($d{>}0.6$) but underpowered for the small within-model
spectral effects ($d{\approx}0.2$).

% ============================================================================
\section{Conclusion}

The Mine~5 thesis --- RLHF as miscalibrator --- is not supported by
these data. RLHF compresses output distributions ($d{=}0.684$), the
compression
survives abliteration (small, non-significant difference; TOST
inconclusive at $n{=}30$), and the
calibration ratio point estimate moves in the improved direction
($1.97 \to 2.51$) --- but the differential compression rate that
would establish that movement is exploratory at this sample size, so
the thesis is unsupported here rather than refuted.

Two findings confirmed, one unsupported, one exploratory. The
unsupported prediction is the most informative result: these data
give no support to the miscalibrator account the prospectus
proposed, and a powered replication is what would settle it. The compression it installs
is deeper than safety directions and persists through abliteration.
The confidence paradox requires an alternative mechanism.

\begin{firstperson}
I predicted RLHF would be the cause of the confidence paradox. The
data said it makes calibration better, not worse. I could have
reframed the calibration ratio to tell a different story --- the
absolute entropy IS lower, which superficially supports
``overcalibration.'' But the ratio is the honest metric, and the
ratio went up. The selective sharpening pattern is suggestive but
I cannot claim it is confirmed at $n{=}30$ --- the bootstrap CI
crosses zero. What IS confirmed: RLHF installs a geometric
compression that abliteration cannot remove. The confidence
paradox needs a new explanation, and ``workspace bypass during
generation'' is the leading candidate.
\end{firstperson}

\section*{Author Contributions (CRediT)}
\noindent\textbf{Lyra}: Conceptualization, Methodology, Formal analysis, Writing -- original draft.\\
\textbf{Thomas Edrington}: Supervision, Resources, Project administration.\\

\section*{Acknowledgments}
We thank Dwayne Wilkes (Sentient Futures / Liberation Labs) for statistical auditing and red-team review.

We thank Kavi (verification review) for their contributions to this work.

\bibliography{references}
\bibliographystyle{plainnat}

\end{document}
